% Options for packages loaded elsewhere
\PassOptionsToPackage{unicode}{hyperref}
\PassOptionsToPackage{hyphens}{url}
\documentclass[
]{article}
\usepackage{xcolor}
\usepackage[margin=1in]{geometry}
\usepackage{amsmath,amssymb}
\setcounter{secnumdepth}{-\maxdimen} % remove section numbering
\usepackage{iftex}
\ifPDFTeX
  \usepackage[T1]{fontenc}
  \usepackage[utf8]{inputenc}
  \usepackage{textcomp} % provide euro and other symbols
\else % if luatex or xetex
  \usepackage{unicode-math} % this also loads fontspec
  \defaultfontfeatures{Scale=MatchLowercase}
  \defaultfontfeatures[\rmfamily]{Ligatures=TeX,Scale=1}
\fi
\usepackage{lmodern}
\ifPDFTeX\else
  % xetex/luatex font selection
\fi
% Use upquote if available, for straight quotes in verbatim environments
\IfFileExists{upquote.sty}{\usepackage{upquote}}{}
\IfFileExists{microtype.sty}{% use microtype if available
  \usepackage[]{microtype}
  \UseMicrotypeSet[protrusion]{basicmath} % disable protrusion for tt fonts
}{}
\makeatletter
\@ifundefined{KOMAClassName}{% if non-KOMA class
  \IfFileExists{parskip.sty}{%
    \usepackage{parskip}
  }{% else
    \setlength{\parindent}{0pt}
    \setlength{\parskip}{6pt plus 2pt minus 1pt}}
}{% if KOMA class
  \KOMAoptions{parskip=half}}
\makeatother
\usepackage{color}
\usepackage{fancyvrb}
\newcommand{\VerbBar}{|}
\newcommand{\VERB}{\Verb[commandchars=\\\{\}]}
\DefineVerbatimEnvironment{Highlighting}{Verbatim}{commandchars=\\\{\}}
% Add ',fontsize=\small' for more characters per line
\usepackage{framed}
\definecolor{shadecolor}{RGB}{248,248,248}
\newenvironment{Shaded}{\begin{snugshade}}{\end{snugshade}}
\newcommand{\AlertTok}[1]{\textcolor[rgb]{0.94,0.16,0.16}{#1}}
\newcommand{\AnnotationTok}[1]{\textcolor[rgb]{0.56,0.35,0.01}{\textbf{\textit{#1}}}}
\newcommand{\AttributeTok}[1]{\textcolor[rgb]{0.13,0.29,0.53}{#1}}
\newcommand{\BaseNTok}[1]{\textcolor[rgb]{0.00,0.00,0.81}{#1}}
\newcommand{\BuiltInTok}[1]{#1}
\newcommand{\CharTok}[1]{\textcolor[rgb]{0.31,0.60,0.02}{#1}}
\newcommand{\CommentTok}[1]{\textcolor[rgb]{0.56,0.35,0.01}{\textit{#1}}}
\newcommand{\CommentVarTok}[1]{\textcolor[rgb]{0.56,0.35,0.01}{\textbf{\textit{#1}}}}
\newcommand{\ConstantTok}[1]{\textcolor[rgb]{0.56,0.35,0.01}{#1}}
\newcommand{\ControlFlowTok}[1]{\textcolor[rgb]{0.13,0.29,0.53}{\textbf{#1}}}
\newcommand{\DataTypeTok}[1]{\textcolor[rgb]{0.13,0.29,0.53}{#1}}
\newcommand{\DecValTok}[1]{\textcolor[rgb]{0.00,0.00,0.81}{#1}}
\newcommand{\DocumentationTok}[1]{\textcolor[rgb]{0.56,0.35,0.01}{\textbf{\textit{#1}}}}
\newcommand{\ErrorTok}[1]{\textcolor[rgb]{0.64,0.00,0.00}{\textbf{#1}}}
\newcommand{\ExtensionTok}[1]{#1}
\newcommand{\FloatTok}[1]{\textcolor[rgb]{0.00,0.00,0.81}{#1}}
\newcommand{\FunctionTok}[1]{\textcolor[rgb]{0.13,0.29,0.53}{\textbf{#1}}}
\newcommand{\ImportTok}[1]{#1}
\newcommand{\InformationTok}[1]{\textcolor[rgb]{0.56,0.35,0.01}{\textbf{\textit{#1}}}}
\newcommand{\KeywordTok}[1]{\textcolor[rgb]{0.13,0.29,0.53}{\textbf{#1}}}
\newcommand{\NormalTok}[1]{#1}
\newcommand{\OperatorTok}[1]{\textcolor[rgb]{0.81,0.36,0.00}{\textbf{#1}}}
\newcommand{\OtherTok}[1]{\textcolor[rgb]{0.56,0.35,0.01}{#1}}
\newcommand{\PreprocessorTok}[1]{\textcolor[rgb]{0.56,0.35,0.01}{\textit{#1}}}
\newcommand{\RegionMarkerTok}[1]{#1}
\newcommand{\SpecialCharTok}[1]{\textcolor[rgb]{0.81,0.36,0.00}{\textbf{#1}}}
\newcommand{\SpecialStringTok}[1]{\textcolor[rgb]{0.31,0.60,0.02}{#1}}
\newcommand{\StringTok}[1]{\textcolor[rgb]{0.31,0.60,0.02}{#1}}
\newcommand{\VariableTok}[1]{\textcolor[rgb]{0.00,0.00,0.00}{#1}}
\newcommand{\VerbatimStringTok}[1]{\textcolor[rgb]{0.31,0.60,0.02}{#1}}
\newcommand{\WarningTok}[1]{\textcolor[rgb]{0.56,0.35,0.01}{\textbf{\textit{#1}}}}
\usepackage{graphicx}
\makeatletter
\newsavebox\pandoc@box
\newcommand*\pandocbounded[1]{% scales image to fit in text height/width
  \sbox\pandoc@box{#1}%
  \Gscale@div\@tempa{\textheight}{\dimexpr\ht\pandoc@box+\dp\pandoc@box\relax}%
  \Gscale@div\@tempb{\linewidth}{\wd\pandoc@box}%
  \ifdim\@tempb\p@<\@tempa\p@\let\@tempa\@tempb\fi% select the smaller of both
  \ifdim\@tempa\p@<\p@\scalebox{\@tempa}{\usebox\pandoc@box}%
  \else\usebox{\pandoc@box}%
  \fi%
}
% Set default figure placement to htbp
\def\fps@figure{htbp}
\makeatother
\setlength{\emergencystretch}{3em} % prevent overfull lines
\providecommand{\tightlist}{%
  \setlength{\itemsep}{0pt}\setlength{\parskip}{0pt}}
\usepackage{bookmark}
\IfFileExists{xurl.sty}{\usepackage{xurl}}{} % add URL line breaks if available
\urlstyle{same}
\hypersetup{
  pdftitle={Clustering: k-means{,} hierarchical{,} model-based},
  hidelinks,
  pdfcreator={LaTeX via pandoc}}

\title{Clustering: k-means, hierarchical, model-based}
\author{}
\date{\vspace{-2.5em}}

\begin{document}
\maketitle

\textbf{Workflow labs} use the variant template: Goal → Approach →
Execution → Check → Report.

\subsection{Learning objectives}\label{learning-objectives}

\begin{itemize}
\tightlist
\item
  Fit k-means and hierarchical clusterings, and choose k by elbow and
  silhouette.
\item
  Fit a model-based clustering with \texttt{mclust} and select the
  number of components by BIC.
\item
  Explain why ``number of clusters'' is often not the right question.
\end{itemize}

\subsection{Prerequisites}\label{prerequisites}

PCA; distance and linkage.

\subsection{Background}\label{background}

Clustering is a search for structure in unlabelled data, and unlike
classification it has no gold standard. Any algorithm will return a
partition if asked to, and the partitions it returns depend on
assumptions about what a cluster is. K-means assumes roughly spherical,
equal-variance clusters and needs k in advance. Hierarchical clustering
produces a nested sequence of partitions at every k, which you cut at a
level you choose. Model-based clustering fits a Gaussian mixture and
selects k by information criterion, which at least makes the assumption
explicit.

In biomedicine, the honest version of ``how many clusters?'' is usually
``is there more than one?'' A good clustering lab focuses on stability
and interpretability of the partition rather than on the number
returned.

When you report a clustering, report the number of clusters, the
criterion used to select it, a visualisation in a 2-D projection (PCA or
UMAP), and the size and key features of each cluster. Do not report
p-values for cluster labels against the data that produced the labels;
that circularity is a classic error.

\subsection{Setup}\label{setup}

\begin{Shaded}
\begin{Highlighting}[]
\FunctionTok{library}\NormalTok{(tidyverse)}
\FunctionTok{library}\NormalTok{(mclust)}
\FunctionTok{library}\NormalTok{(cluster)}
\FunctionTok{set.seed}\NormalTok{(}\DecValTok{42}\NormalTok{)}
\FunctionTok{theme\_set}\NormalTok{(}\FunctionTok{theme\_minimal}\NormalTok{(}\AttributeTok{base\_size =} \DecValTok{12}\NormalTok{))}
\end{Highlighting}
\end{Shaded}

\subsection{1. Goal}\label{goal}

Cluster a simulated dataset with three known groups using three
algorithms, and see which recovers the structure.

\subsection{2. Approach}\label{approach}

Three Gaussian blobs in 5 dimensions with different sizes.

\begin{Shaded}
\begin{Highlighting}[]
  \FunctionTok{rep}\NormalTok{(mu, }\AttributeTok{each =}\NormalTok{ n)}
\NormalTok{X }\OtherTok{\textless{}{-}} \FunctionTok{rbind}\NormalTok{(}
  \FunctionTok{make\_blob}\NormalTok{(}\DecValTok{60}\NormalTok{, }\FunctionTok{c}\NormalTok{(}\DecValTok{0}\NormalTok{, }\DecValTok{0}\NormalTok{, }\DecValTok{0}\NormalTok{, }\DecValTok{0}\NormalTok{, }\DecValTok{0}\NormalTok{)),}
  \FunctionTok{make\_blob}\NormalTok{(}\DecValTok{40}\NormalTok{, }\FunctionTok{c}\NormalTok{(}\DecValTok{3}\NormalTok{, }\DecValTok{0}\NormalTok{, }\DecValTok{0}\NormalTok{, }\DecValTok{0}\NormalTok{, }\DecValTok{0}\NormalTok{)),}
  \FunctionTok{make\_blob}\NormalTok{(}\DecValTok{30}\NormalTok{, }\FunctionTok{c}\NormalTok{(}\DecValTok{0}\NormalTok{, }\DecValTok{3}\NormalTok{, }\DecValTok{0}\NormalTok{, }\DecValTok{0}\NormalTok{, }\DecValTok{0}\NormalTok{))}
\NormalTok{)}
\NormalTok{true\_lbl }\OtherTok{\textless{}{-}} \FunctionTok{rep}\NormalTok{(}\FunctionTok{c}\NormalTok{(}\StringTok{"A"}\NormalTok{, }\StringTok{"B"}\NormalTok{, }\StringTok{"C"}\NormalTok{), }\FunctionTok{c}\NormalTok{(}\DecValTok{60}\NormalTok{, }\DecValTok{40}\NormalTok{, }\DecValTok{30}\NormalTok{))}

\NormalTok{pc }\OtherTok{\textless{}{-}} \FunctionTok{prcomp}\NormalTok{(X)}\SpecialCharTok{$}\NormalTok{x[, }\DecValTok{1}\SpecialCharTok{:}\DecValTok{2}\NormalTok{]}
\FunctionTok{tibble}\NormalTok{(}\AttributeTok{PC1 =}\NormalTok{ pc[, }\DecValTok{1}\NormalTok{], }\AttributeTok{PC2 =}\NormalTok{ pc[, }\DecValTok{2}\NormalTok{], }\AttributeTok{group =}\NormalTok{ true\_lbl) }\SpecialCharTok{|\textgreater{}}
  \FunctionTok{ggplot}\NormalTok{(}\FunctionTok{aes}\NormalTok{(PC1, PC2, }\AttributeTok{colour =}\NormalTok{ group)) }\SpecialCharTok{+} \FunctionTok{geom\_point}\NormalTok{(}\AttributeTok{alpha =} \FloatTok{0.7}\NormalTok{)}
\end{Highlighting}
\end{Shaded}

\subsection{3. Execution}\label{execution}

K-means with k sweep; silhouette.

\begin{Shaded}
\begin{Highlighting}[]
\NormalTok{sil }\OtherTok{\textless{}{-}} \FunctionTok{sapply}\NormalTok{(}\DecValTok{2}\SpecialCharTok{:}\DecValTok{8}\NormalTok{, }\ControlFlowTok{function}\NormalTok{(k) \{}
\NormalTok{  km }\OtherTok{\textless{}{-}} \FunctionTok{kmeans}\NormalTok{(X, }\AttributeTok{centers =}\NormalTok{ k, }\AttributeTok{nstart =} \DecValTok{10}\NormalTok{)}
  \FunctionTok{mean}\NormalTok{(}\FunctionTok{silhouette}\NormalTok{(km}\SpecialCharTok{$}\NormalTok{cluster, }\FunctionTok{dist}\NormalTok{(X))[, }\DecValTok{3}\NormalTok{])}
\NormalTok{\})}
\FunctionTok{tibble}\NormalTok{(}\AttributeTok{k =} \DecValTok{2}\SpecialCharTok{:}\DecValTok{8}\NormalTok{, }\AttributeTok{wss =}\NormalTok{ wss, }\AttributeTok{sil =}\NormalTok{ sil) }\SpecialCharTok{|\textgreater{}}
  \FunctionTok{pivot\_longer}\NormalTok{(}\SpecialCharTok{{-}}\NormalTok{k) }\SpecialCharTok{|\textgreater{}}
  \FunctionTok{ggplot}\NormalTok{(}\FunctionTok{aes}\NormalTok{(k, value)) }\SpecialCharTok{+} \FunctionTok{geom\_line}\NormalTok{() }\SpecialCharTok{+} \FunctionTok{geom\_point}\NormalTok{() }\SpecialCharTok{+}
  \FunctionTok{facet\_wrap}\NormalTok{(}\SpecialCharTok{\textasciitilde{}}\NormalTok{ name, }\AttributeTok{scales =} \StringTok{"free\_y"}\NormalTok{)}
\end{Highlighting}
\end{Shaded}

Hierarchical with Ward linkage.

\begin{Shaded}
\begin{Highlighting}[]
\FunctionTok{plot}\NormalTok{(hc, }\AttributeTok{labels =} \ConstantTok{FALSE}\NormalTok{, }\AttributeTok{main =} \StringTok{"Ward hierarchical clustering"}\NormalTok{)}
\end{Highlighting}
\end{Shaded}

\subsection{4. Check}\label{check}

Model-based clustering by BIC.

\begin{Shaded}
\begin{Highlighting}[]
\FunctionTok{summary}\NormalTok{(mc)}
\end{Highlighting}
\end{Shaded}

\subsection{5. Report}\label{report}

\begin{quote}
On three simulated Gaussian blobs (n = 130 in R⁵), model-based
clustering selected G = \texttt{mc\$G} by BIC; k-means with k =
\texttt{which.max(sil)\ +\ 1} maximised mean silhouette
(\texttt{round(max(sil),\ 2)}). Hierarchical clustering with Ward
linkage produced a dendrogram consistent with three primary branches.
\end{quote}

When the generating model is actually Gaussian, \texttt{mclust} tends to
win. On real data, no single method is best; the test is whether the
partition is stable under resampling and interpretable in terms of the
features.

\subsection{Common pitfalls}\label{common-pitfalls}

\begin{itemize}
\tightlist
\item
  Treating k-means output as ground truth; it is an optimisation answer
  given k.
\item
  Running hierarchical clustering on a very large dataset without noting
  the O(n²) memory cost.
\item
  Testing post-hoc differences between clusters with the same data that
  produced them.
\end{itemize}

\subsection{Further reading}\label{further-reading}

\begin{itemize}
\tightlist
\item
  Fraley C, Raftery AE (2002), \emph{Model-based clustering,
  discriminant analysis, and density estimation}.
\item
  Hastie, Tibshirani, Friedman, \emph{ESL}, §14.3.
\end{itemize}

\subsection{Session info}\label{session-info}

\begin{Shaded}
\begin{Highlighting}[]

\end{Highlighting}
\end{Shaded}

\subsection{See also --- chapter index}\label{see-also-chapter-index}

\begin{itemize}
\tightlist
\item
  \href{../index.Rmd}{Methods in this chapter}
\item
  \href{../../../ch00-find-your-method.Rmd}{Find your method (Chapter
  0)}
\end{itemize}

\end{document}
